% Abhakseite „Das kann ich“ – Zone nur im Gesamt (\mitzone)
%% TODO Ich-kann-Sätze: Platzhalter wie in den Aufgabendateien
\begin{abhakseite}
\ifdefined\mitzone
\abhakgruppe{Kennst du schon}
\abhak{1}{Geradengleichungen aufstellen und den allgemeinen Geradenpunkt bilden}
\abhak{2}{Ebenengleichungen lesen (Koordinatenform, Parameterform, Normalenvektor)}
\abhak{3}{Den Lagebefund vorweg führen (liegt der Punkt darin, schneidet die Gerade überhaupt – Merkregel über das Skalarprodukt)}
\abhak{4}{Lineare Gleichungssysteme mit zwei und drei Unbekannten lösen, auch überbestimmte mit Probegleichung}
\abhak{5}{Teilverhältnisse und Parameterbereiche von Strecken lesen}
\abhak{6}{Schrägbilder lesen und zeichnen}
\abhak{7}{Geradengleichungen aufstellen und den allgemeinen Geradenpunkt bilden – Fehler finden}
\abhak{8}{Geradengleichungen aufstellen und den allgemeinen Geradenpunkt bilden – selbst rechnen}
\fi
\abhakgruppe{Einheit 1 · Schnittpunkt von Gerade und Ebene}
\abhak{9}{Was wird geschnitten – und was kann herauskommen?}
\abhak{10}{Schnittpunkt von Gerade und Ebene}
\abhak{11}{Höhe einer Spitze aus dem Abstand ihres Schattenpunkts zu einem Punkt bestimmen}
\abhak{12}{Schnittpunkt von Gerade und Ebene – Fehler finden, Begründen}
\abhakgruppe{Einheit 2 · Schnittpunkt zweier Geraden}
\abhak{13}{Schnittpunkt zweier Geraden}
\abhak{14}{Schattenpunkt auf einer Kante als Schnittpunkt von Lichtgerade und Kantengerade berechnen}
\abhak{15}{Schnittpunkt zweier Geraden – Fehler finden, Begründen}
\abhakgruppe{Einheit 3 · Spuren, Schnittgeraden und Schnittfiguren}
\abhak{16}{Spuren, Schnittgeraden und Schnittfiguren}
\abhak{17}{Spuren, Schnittgeraden und Schnittfiguren – Fehler finden, Begründen}
\end{abhakseite}
